\documentclass[10pt,a4paper]{amsart}
\usepackage[margin=19mm]{geometry}
\usepackage{amsmath,amssymb,mathtools}
\usepackage[scr=rsfs,cal=euler]{mathalfa}
\usepackage{xfrac}
\usepackage[unicode,hidelinks,bookmarks=false]{hyperref}
\newcommand{\ie}{i.e.\ }
\newcommand{\cf}{cf.\ }
\newcommand{\resp}{respectively\ }
\newcommand{\Subsection}{\subsection}
\providecommand{\nobreakdash}{\nobreak\hbox{-}}
\setlength{\emergencystretch}{2em}
\multlinegap0pt
\usepackage{amssymb}
\usepackage{tikz}
\usepackage{bm}
\usepackage{mathtools}
\usepackage{xcolor}
\newtheorem{lemma}{Lemma}
\title{Augmentation and Bulk Edge Correspondence for one dimensional aperiodic tight binding operators}
\author{Johannes Kellendonk, Lorenzo Scaglione}
\date{Source: arXiv:2512.01555v2 (CC BY 4.0)}
\begin{document}
\maketitle

\noindent\emph{Source and licence.} Augmentation and Bulk Edge Correspondence for one dimensional aperiodic tight binding operators. Version arXiv:2512.01555v2 (CC BY 4.0), distributed by arXiv under the Creative Commons Attribution 4.0 International licence, http://creativecommons.org/licenses/by/4.0/, as recorded in the arXiv licence metadata for this exact version and shown on the abstract page https://arxiv.org/abs/2512.01555v2. Full licence notice: \texttt{licenses/arxiv-2512.01555-CC-BY-4.0.txt}.\par\smallskip

\noindent\textit{Editorial scope.} Counted unit: the article's lemma lem-tr-ch (source0.tex lines 1520--1527). The article's complete original proof of it is delivered with it (source0.tex lines 1529--1571, one \texttt{proof} environment of 43 lines). No mathematics is written, repaired or re-derived: the statement and the proof are the article's own lines, in the article's own notation, and no retained line is substituted or re-flowed.  No further source-local context is needed: every cross-reference of every retained line resolves inside the two delivered spans.  Nothing was omitted from any retained span, so the package carries no omission record.  No retained line contains a citation, so the wrapper carries no bibliography and no bibliography entry.  No retained line contains a figure environment, an image-inclusion command or drawing code, so no image is delivered and the package declares no asset dependency.  Editorial formatting only, in addition: an AMS \texttt{amsart} preamble that restates the article's own declarations and macros used by the retained lines, namely the environments \texttt{lemma} on separate counters, together with the standard packages those lines need.  The retained environments print in this document as Lemma 1 in order of appearance, whereas the article prints the same items with the numbering of its own full document.  The complete native source of the article as distributed in the arXiv e-print for this exact version is bundled unchanged as \texttt{sources/190130/source0.tex}.
\begin{lemma}\label{lem-tr-ch}
Given a projection $p\in A \rtimes_\alpha \mathbb{Z}$ let $P$ be as above.
Then 
\begin{equation}\label{eq-tr-ch}
\frac1{\mathrm{i}}\hat{\tau}_{\tilde{A}}(P[\delta_1(P), \delta_2(P)])
=\hat{\tau}_A(p).
\end{equation}
\end{lemma}

\begin{proof} 
Clearly $P(0)=\begin{pmatrix}
		0 & 0\\
		0 & p
	\end{pmatrix}$ and $P(1)=\begin{pmatrix}
		0 & 0\\
		0 & \bar{\alpha}(p)
	\end{pmatrix}$ so that $P$ is a projection in $M_2(M_{\bar{\alpha}}(A \rtimes_\alpha \mathbb{Z}))$ and $q(P)=0 \oplus p$. 
	Denote $c_t=\cos{(\frac{\pi}{2}t)}$ and $s_t=\sin{(\frac{\pi}{2}t)}$. A direct calculation leads to (we write shorter $\dot f=\delta_1(f)$ and $f' = \delta_2(f)$): 
	\begin{align*} 
		\dot{P}_t &=  \frac{\pi}{2}U_t (A_1+A_2)U_t^*,\quad A_1=p\begin{pmatrix}
			0 & 1\\ 
			1 & 0
		\end{pmatrix}, \quad A_2=\begin{pmatrix}
			0 & -(c_t^2u^*+s_t^2u)p\\ 
			-p(c_t^2u+s_t^2u^*) & c_t s_t[p,u^*-u]
		\end{pmatrix}\\
		P'_t &= U_t (B_1+B_2)U_t^*,\quad B_1=p'\begin{pmatrix}
			0 & 0\\ 
			0 & 1
		\end{pmatrix}, \quad B_2=\mathrm{i}c_t s_t p\begin{pmatrix}
			0 & 1\\ 
			-1 & 0
		\end{pmatrix}.
	\end{align*}
	Hence
	$$\frac{1}{\mathrm{i}}\hat{\tau}_{\tilde{A}}(P[\dot{P}, P'])=\frac{\pi}{2\mathrm{i}}\int_0^1\hat{\tau}_A\left( U_t \begin{pmatrix}
		0 & 0\\ 0 & p
	\end{pmatrix} ([A_1,B_1]+ [A_1,B_2]+[A_2,B_1]+[A_2,B_2])U_t^*\right)\mathrm{d}t.$$
	We compute the four terms. We see easily that the first term is zero and second one is equal to $\hat{\tau}_A(p)$. For the third term, we get
	\begin{align*}
		\mathrm{3\textsuperscript{rd} term}&= \frac{\pi}{2\mathrm{i}}\int_0^1\hat{\tau}_A\left( \begin{pmatrix}
			0 & 0\\ 0 & p
		\end{pmatrix} [A_2,B_1]\right)\mathrm{d}t\\ &= 
		\frac{1}{2\mathrm{i}}\hat{\tau}_A(p(u^*-u)p'-p(u^*-u)pp'-pp'p(u^*-u)+p'(u^*-u)p)\\
		&= \frac{1}{2\mathrm{i}}\hat{\tau}_A((p(u^*-u)p)'-p(u^*-u)'p)\\
		&= \frac{1}{2}\hat{\tau}_A(p(u^*+u)p)),
	\end{align*}
	where we used the relation $pp'p=0$, the cyclicity of the trace and that it vanishes on a total derivative. 
	Finally
	$$\mathrm{4\textsuperscript{th} term}=-\frac{1}{2}\hat{\tau}_A(p(u^*+u)p)$$
	follows quickly. Thus only the second term contributes yielding the result.
\end{proof}

\end{document}
